% ============================================================
% content/solutions.tex —— 2026 湖南卷物理 · 参考答案与详细解析
% 解析类型：source_embedded（转录自官方 Word 版内嵌解析，独立复核）
% ============================================================

\begin{answer}
\gksolhead{1}{单项选择题}
\gkanswer{答案 B}
由加速度的定义式
\[
  a=\frac{\Delta v}{\Delta t}=\frac{50-30}{1}\,\si{\metre\per\second\squared}
   =\SI{20}{\metre\per\second\squared},
\]
故选 B。
\end{answer}

\begin{answer}
\gksolhead{2}{单项选择题}
\gkanswer{答案 D}
$\alpha$ 衰变放出的 $\alpha$ 粒子为 ${}_{2}^{4}\mathrm{He}$，衰变过程电荷数守恒，因此母核电荷数减少 2、质量数减少 4，故 A、B 错误；$\alpha$ 粒子电离能力强、穿透能力极弱，一张白纸即可将其挡住；穿透能力强、可穿透几厘米厚铅板的是 $\gamma$ 射线，故 C 错误；衰变方程满足电荷数守恒和质量数守恒：左边电荷数为 86，右边 $2\times2+82=86$；左边质量数为 216，右边 $2\times4+208=216$，守恒关系成立，故 D 正确。
\end{answer}

\begin{answer}
\gksolhead{3}{单项选择题}
\gkanswer{答案 C}
设恒星半径为 $R$，由题意得行星轨道半径 $r=nR$。行星绕恒星做匀速圆周运动，万有引力提供向心力，有
\[
  G\frac{Mm}{r^{2}}=m\frac{4\pi^{2}}{T^{2}}\,r .
\]
整理得恒星质量 $M=\dfrac{4\pi^{2}r^{3}}{GT^{2}}$；恒星为均匀球体，体积 $V=\dfrac{4}{3}\pi R^{3}$，平均密度 $\rho=\dfrac{M}{V}$，将 $r=nR$ 代入得
\[
  \rho=\frac{\frac{4\pi^{2}(nR)^{3}}{GT^{2}}}{\frac{4}{3}\pi R^{3}}
      =\frac{3\pi n^{3}}{GT^{2}},
\]
故选 C。
\end{answer}

\begin{answer}
\gksolhead{4}{单项选择题}
\gkanswer{答案 B}
某时刻 $A$ 处质点位于波峰、$B$ 处质点位于波谷，故 $A$、$B$ 两点的波程差应满足半波长的奇数倍
\[
  \Delta r=\left(n+\frac{1}{2}\right)\lambda,\qquad n=0,1,2,\dots
\]
若波源在 $(0,0)$，$A$ 点到波源距离
\[
  r_{A}=\sqrt{4^{2}+3^{2}}\,\si{\metre}=5\,\si{\metre},\qquad r_{B}=4\,\si{\metre},
\]
波程差 $\Delta r=1\,\si{\metre}$，代入得 $1=\left(n+\tfrac{1}{2}\right)\lambda$，解得 $\lambda=\dfrac{2}{2n+1}\,\si{\metre}$，故 B 正确；若波源在 $(0,3)$，波程差仍为 $1\,\si{\metre}$，故 C、D 错误。
\end{answer}

\begin{answer}
\gksolhead{5}{单项选择题}
\gkanswer{答案 A}
设单个电池组的电动势为 $E$，内阻为 $r$。$a$、$b$ 接 $d$、$e$ 时只有单个电池接入电路，由闭合电路欧姆定律：
\[
  E=I_{1}(R+r).
\]
$a$、$c$ 接 $d$、$e$ 时两个电池串联，总电动势为 $2E$、总内阻为 $2r$，同理得
\[
  2E=I_{2}(R+2r).
\]
代入 $R=\SI{10}{\ohm}$、$I_{1}=\SI{1.0}{\ampere}$、$I_{2}=\SI{1.5}{\ampere}$ 联立解得 $E=\SI{15}{\volt}$，$r=\SI{5}{\ohm}$，故选 A。
\end{answer}

\begin{answer}
\gksolhead{6}{单项选择题}
\gkanswer{答案 A}
$A$、$C$ 板均接地，电势为 0，$B$ 板是等势体，故 $BA$ 间、$BC$ 间的电势差相等，都等于 $U=\varphi_{B}$。平行板内部为匀强电场，场强满足 $E=\dfrac{U}{d}$，已知 $d_{1}<d_{2}$，故 $E_{1}>E_{2}$，A 正确、B 错误。$a$ 点到 $A$ 板的距离为 $\tfrac{d_{1}}{2}$，电势
\[
  \varphi_{a}=E_{1}\cdot\frac{d_{1}}{2}=\frac{U}{d_{1}}\cdot\frac{d_{1}}{2}=\frac{U}{2},
\]
同理 $\varphi_{b}=\dfrac{U}{2}$，故 $\varphi_{a}=\varphi_{b}$，C、D 错误。
\end{answer}

\begin{answer}
\gksolhead{7}{多项选择题}
\gkanswer{答案 D}
设心形气球总体积为 $V$（恒定），小气球体积为 $V_{1}$，则外部气体体积 $V_{2}=V-V_{1}$。设小气球内气体压强 $p_{1}$、外部气体压强 $p_{2}$，且 $p_{1}>p_{2}$。结合克拉珀龙方程 $pV=nRT$ 可知 $n_{1}>n_{2}$，即小气球内部气体分子数较多，故 B 错误。

温度升高时，假设小气球内、外部气体体积均不变，由查理定律 $\dfrac{\Delta p}{\Delta T}=\dfrac{p}{T}$，解得小气球内、外气体增加的压强 $\Delta p=\dfrac{\Delta T}{T}\,p$。因 $p_{1}>p_{2}$，升高相同的温度，小气球内部气体压强增加得比外部多，故小气球内部气体的体积将变大、外部气体的体积将变小，故 A、C 错误，D 正确。
\end{answer}

\begin{answer}
\gksolhead{8}{多项选择题}
\gkanswer{答案 BC}
温度是分子平均动能的统计标志，仅反映整体的平均水平：高温处仍存在速率小的水分子，低温处也存在速率大的水分子，并非所有分子的运动速率都满足高温处大于低温处，故 A 错误。一定量的水变为同温度的水蒸气，温度不变则分子总动能不变；汽化过程吸收热量，且汽化后分子间距大幅增大、分子势能增加，因此总内能变大，故 B 正确。该过程中，均热板从高温的芯片吸热、向低温的冷凝端放热，整个过程循环进行，均热板自身状态复原，不需要对外界做功，符合物理规律，故 C 正确。根据热力学第二定律，不可能从低温物体吸热、向高温物体放热，同时不产生其他影响，D 错误。
\end{answer}

\begin{answer}
\gksolhead{9}{多项选择题}
\gkanswer{答案 BC}
设每根绳子的拉力为 $F_{T}$，两根绳子之间的夹角为 $2\alpha$，每个气球受到空气阻力 $f=kv$、浮力为 $F_{\text{浮}}$。对整体受力分析，由平衡条件，竖直方向上有
\[
  2F_{\text{浮}}=2G+2F_{T}\cos\alpha\sin\theta,
\]
水平方向上有
\[
  2kv=2F_{T}\cos\alpha\cos\theta,
\]
整理可得 $\tan\theta=\dfrac{F_{\text{浮}}-G}{kv}$。运动速度缓慢增大，$\tan\theta$ 减小，则 $\theta$ 减小、$F_{T}$ 增大，故 A 错误、B 正确。对单个气球受力分析，由平衡条件有 $F_{T}\sin\alpha=F$，由于 $\alpha$ 不变、$F_{T}$ 增大，则 $F$ 增大，故 C 正确、D 错误。
\end{answer}

\begin{answer}
\gksolhead{10}{多项选择题}
\gkanswer{答案 AD}
小球所受洛伦兹力的水平分力提供小球做圆周运动的向心力，由左手定则可知，从管道上方俯视，小球沿顺时针方向做圆周运动，故 A 正确。

设磁感应强度为 $B$、小球的速度为 $v$，此时小球与管道间无弹力。竖直方向上有
\[
  mg=qvB\cos45^{\circ},
\]
水平方向上有
\[
  qvB\sin45^{\circ}=m\frac{v^{2}}{R},
\]
解得 $v=\sqrt{gR}$，则小球做圆周运动的周期
\[
  T=\frac{2\pi R}{v}=2\pi\sqrt{\frac{R}{g}},
\]
故 B 错误。水平方向上，由动量定理
\[
  I_{x}=\Delta p=-mv-mv=-2m\sqrt{gR};
\]
竖直方向有
\[
  I_{y}=qvB\cos45^{\circ}\cdot\frac{T}{2}=mg\cdot\frac{T}{2}=m\pi\sqrt{gR},
\]
则半个周期内洛伦兹力的冲量大小
\[
  I=\sqrt{I_{x}^{2}+I_{y}^{2}}=m\sqrt{(4+\pi^{2})gR},
\]
故 C 错误。若速率为 2 倍，竖直方向上有 $mg+F_{Ny}=q\cdot 2vB\cos45^{\circ}$，水平方向上有 $q\cdot 2vB\sin45^{\circ}+F_{Nx}=m\dfrac{(2v)^{2}}{R}$，弹力大小 $F_{N}=\sqrt{F_{Ny}^{2}+F_{Nx}^{2}}$，联立解得 $F_{N}=\sqrt{5}\,mg$，故 D 正确。
\end{answer}

\begin{answer}
\gksolhead{11}{实验题}
\gkanswer{答案 (3) $0.6\pi$；(4) $b$；(5) $g=\dfrac{4\pi^{2}l}{T^{2}}$；(6) 高于}
(3) 由题图 1 可知，单摆的周期 $T=0.6\pi\,\si{\second}$。

(4) 用刻度尺测量摆长时，应读出球心所在位置的读数，即 $b$ 位置的读数，该读数即为摆长 $l$。

(5) 由单摆周期公式 $T=2\pi\sqrt{\dfrac{l}{g}}$ 可得重力加速度
\[
  g=\frac{4\pi^{2}l}{T^{2}}.
\]

(6) 若所得实验结果小于当地重力加速度，可知摆长测量值偏小，即刻度尺 0 刻度线高于悬点。
\end{answer}

\begin{answer}
\gksolhead{12}{实验题}
\gkanswer{答案 (2) $c$、$d$；50；(3) 11.350；(5) 见解析；(6) 1.55}
(2) 由霍尔效应原理可知，霍尔电压出现在垂直于电流和磁场的方向上，结合图 2 可知，应将电压表接在霍尔元件的 $c$、$d$ 两端。由图 3 可知，电压表的最小刻度为 \SI{2}{\milli\volt}，则此时电压为 \SI{50}{\milli\volt}。

(3) 图 4 中螺旋测微器读数为 $11\,\si{\milli\metre}+35.0\times0.01\,\si{\milli\metre}=\SI{11.350}{\milli\metre}$。

(5) 根据表格中数据，在图 5 中描出第 4、5 组测量数据的坐标点，并作出 $z$--$U$ 图像，如图所示。

\gkfig{figures/s12-z-u}

(6) 根据图像可知，电压表示数为 \SI{20}{\milli\volt} 时 $z_{1}=\SI{12.15}{\milli\metre}$，移动后电压表示数为 \SI{60}{\milli\volt} 时 $z_{2}=\SI{13.70}{\milli\metre}$，则此过程中物体位移的大小
\[
  z=z_{2}-z_{1}=\SI{1.55}{\milli\metre}.
\]
\end{answer}

\begin{answer}
\gksolhead{13}{解答题}
\gkanswer{答案 (1) $n=\dfrac{2\sqrt{3}}{3}$；(2) $\sin\theta=0.9$}
(1) 如图所示。

\gkfig{figures/s13-optics}

根据几何关系可知光线 1 的入射角为 $90^{\circ}$，折射角为 $60^{\circ}$，可得
\[
  n=\frac{\sin90^{\circ}}{\sin60^{\circ}}=\frac{2\sqrt{3}}{3}.
\]

(2) 根据几何关系可得光线 2 在介质中的折射角为 $\beta=30^{\circ}$。由折射定律 $n'=\dfrac{\sin\theta}{\sin\beta}$，其中 $n'=1.8$，解得 $\sin\theta=0.9$。
\end{answer}

\begin{answer}
\gksolhead{14}{解答题}
\gkanswer{答案 (1) $a_{0}=\dfrac{BEL}{mr}$；(2) $v=\dfrac{E}{BL}-\dfrac{BLU_{0}}{mk}$；(3) $Q=\dfrac{(R+2r)mE^{2}}{2(R+r)B^{2}L^{2}}$}
(1) $t_{0}$ 时刻开关 $S$ 与 $a$ 连接，流过导体棒的电流
\[
  I=\frac{E}{r},
\]
导体棒所受安培力 $F=BIL$，根据牛顿第二定律 $F=ma_{0}$，联立可得
\[
  a_{0}=\frac{BEL}{mr}.
\]

(2) 当导体棒匀速运动时，导体棒产生的电动势等于电源电动势，设此时导体棒速度为 $v_{0}$，即 $E=BLv_{0}$。当储能器电压为 $U_{0}$ 时，对导体棒根据动量定理有
\[
  -\sum BIL\cdot t=mv-mv_{0},
\]
又 $q=\sum I\cdot t$，可得 $-BLq=mv-mv_{0}$。根据题意 $U=kq$，联立可得
\[
  v=\frac{E}{BL}-\frac{BLU_{0}}{mk}.
\]

(3) 从 $t_{0}\sim t_{1}$ 过程对导体棒，根据动量定理
\[
  B\bar{I}L\cdot t=mv_{0},
\]
又 $q_{1}=\bar{I}\cdot t$，可得 $q_{1}=\dfrac{mE}{B^{2}L^{2}}$。设导体棒上产生的焦耳热为 $Q_{1}$，由能量守恒可得
\[
  q_{1}E=Q_{1}+\frac{1}{2}mv_{0}^{2},
\]
解得
\[
  Q_{1}=\frac{mE^{2}}{2B^{2}L^{2}}.
\]
从 $t_{1}$ 时刻到导体棒静止过程中，根据能量守恒定律可知电路中产生的焦耳热为
\[
  Q_{\text{总}}=\frac{1}{2}mv_{0}^{2}.
\]
导体棒产生的焦耳热为
\[
  Q_{2}=\frac{r}{R+r}\,Q_{\text{总}}
       =\frac{mrE^{2}}{2(R+r)B^{2}L^{2}}.
\]
故
\[
  Q=Q_{1}+Q_{2}
   =\frac{(R+2r)mE^{2}}{2(R+r)B^{2}L^{2}}.
\]
\end{answer}

\begin{answer}
\gksolhead{15}{解答题}
\gkanswer{答案 (1) $v=\sqrt{2gL}$；(2) $\tan\alpha=\dfrac{\sqrt{23}}{2}$}
\gkanswer{(3)(i) $E=\left[\dfrac{(1+k)^{2}}{4k\cos^{2}\theta}-1\right]mgL$；(3)(ii) $d=\dfrac{8}{9}\sqrt{69}\,L\tan\theta$}
(1) 根据机械能守恒定律有
\[
  mgL=\frac{1}{2}mv^{2},
\]
解得小球接触地面瞬间的速度大小
\[
  v=\sqrt{2gL}.
\]

(2) 设小球脱离轻杆时，轻杆与竖直方向的夹角为 $\beta$，如图所示。

\gkfig{figures/s15-rod-beta}

设此时小球的速度为 $v_{1}$，根据机械能守恒定律
\[
  mgL(1-\cos\beta)=\frac{1}{2}mv_{1}^{2},
\]
对小球根据牛顿第二定律（指向圆心方向）
\[
  mg\cos\beta=m\frac{v_{1}^{2}}{L},
\]
解得 $\cos\beta=\dfrac{2}{3}$，$v_{1}=\sqrt{\dfrac{2gL}{3}}$。此时小球水平方向分速度 $v_{x1}=v_{1}\cos\beta$，竖直方向分速度 $v_{y1}=v_{1}\sin\beta$。竖直方向，根据速度位移关系
\[
  v_{y}^{2}-v_{y1}^{2}=2g\cdot L\cos\beta.
\]
小球接触地面瞬间的速度与水平面夹角 $\alpha$ 的正切值 $\tan\alpha=\dfrac{v_{y}}{v_{x1}}$，联立解得
\[
  \tan\alpha=\frac{\sqrt{23}}{2}.
\]

(3)(i) 设其中一部分的质量为 $m_{0}$，则另一部分质量为 $km_{0}$，如图所示。

\gkfig{figures/s15-split}

根据动量守恒定律有
\[
  km_{0}v_{2}\sin\theta=m_{0}v_{3}\sin\theta,\qquad
  km_{0}v_{2}\cos\theta+m_{0}v_{3}\cos\theta=mv.
\]
弹射装置释放的能量
\[
  E=\frac{1}{2}km_{0}v_{2}^{2}+\frac{1}{2}m_{0}v_{3}^{2}-\frac{1}{2}mv^{2},
\]
联立可得
\[
  E=\left[\frac{(1+k)^{2}}{4k\cos^{2}\theta}-1\right]mgL.
\]

(ii) 由均值不等式可知，当 $k=1$ 时 $E$ 最小，为
\[
  E_{\min}=mgL\tan^{2}\theta.
\]
此时 $v_{2}=v_{3}$，两小球同时落地，运动时间为
\[
  t=2\cdot\frac{v_{2}\sin\alpha}{g},
\]
两小球第一次落地点的间距
\[
  d=2v_{2}\sin\theta\cdot t.
\]
由 $\tan\alpha=\dfrac{\sqrt{23}}{2}$ 可得 $\sin\alpha=\sqrt{\dfrac{23}{27}}$，联立解得
\[
  d=\frac{8}{9}\sqrt{69}\,L\tan\theta.
\]
\end{answer}
